\section{Discussion}
\label{sec:discussion}

The measurements describe a quantity that occupies the place of a $\sigma$-profile without being
one. It is sparse where the reference is dense and populated where the reference is empty. It depends on the role the molecule plays in the pair, so the model it feeds does not return the reference state. And it does not decompose, so the familiar move of correcting part of it and expecting part of the benefit is unavailable.

The sharpest statement of the situation is a contrast between two substitutions of the same
values. Supplying the reference profile improves the quantity the profile exists to produce: the
predicted activity coefficients move most of the way toward the measured ones, on every seed. The
same substitution degrades the solubility the pipeline is built to predict. Both facts are measured
on the same checkpoints and the same compounds. Taken together they say that the downstream parts of the model have adapted to the profile the network supplies rather than to the one the physical argument names. Handing the model a physically better value therefore moves it away from the operating point it was fitted at. This is the behaviour the concept-bottleneck literature reports
for learned concepts in general \cite{koh2020cbm,margeloiu2021}, arrived at here through a
thermodynamic decoder and, because the concept has a computable ground truth, measured against the
concept itself rather than inferred from how interventions behave.

It follows that the assumption set out in the Introduction is not a property of physics-informed
models as such but a property of how the internal quantity is trained. The design choice that
matters is whether the gradient of the target reaches the intermediate. Where it does not, as in a
model whose profile predictor is supervised externally and then frozen \cite{yang2025tennet}, the
intermediate should remain what it is named after to within its own fit error, and substitution
should be a sensible operation. That expectation is not tested here, since the model in question
was not retrained. Where it does, the intermediate becomes a free internal coordinate that happens
to be initialised and regularised toward a physical quantity, and the operations the physical
reading licenses stop being safe. The present pipeline sits between the two, with both influences
acting, and the measurements show which one wins.

For practice this suggests a different division of labour than the one the literature has settled
into. If the internal quantity is to be read as chemistry, exchanged with tabulated values, or audited against reference data, it should be supervised against that reference and then held fixed. The capacity needed to fit the data belongs somewhere that makes no physical claim.
If instead the intermediate is allowed to move under the target's gradient, the gain in accuracy
is real but the physical reading has to be dropped, and with it the argument that recommends these
models over direct regression in the first place. What is not available is to train the
intermediate end to end and retain the interpretation, which is the position most of this
literature currently occupies.

Three measurements in the battery are worth separating from the rest because each removes an
explanation that would otherwise be reached for. The output layer is not responsible for the
spurious density: the same head reaches values six orders of magnitude smaller elsewhere in the
same profile. The cost map does not localise the error: the bins the closure is most sensitive to account for no more of the gap than a randomly chosen third. Sensitivity is measured per unit of displaced mass, and most of the mass sits in the insensitive core. And the spurious
density in the empty bins, the most visible defect and the one easiest to describe, is not what the
thermodynamics responds to; the error that propagates is in the magnitude of the profile where the
profile is populated.

The finding is conditional on the fixed relation being an adequate description of the chemistry it
is applied to. Where a relation is well specified for its domain, supplying it with more accurate inputs behaves as the physical reading predicts. A second system, in which a linear free-energy relation stands in for the closure and its adequacy varies by chemical stratum, separates in that way; it is reported in the Supporting Information. The result here should therefore be read as a
statement about an intermediate trained through an approximate closure, which is the common case,
rather than as an argument against building physics into a model.

\section{Limitations}
\label{sec:limitations}

The evidence comes from one pipeline, one closure and one reference tabulation. Every departure reported above is measured against the VT-2005 table. How far that table itself moves with the electronic structure and solvation treatment behind it is a separate question \cite{salih2025sigma}. A different tabulation would shift the reference against which the learned profile is scored, but not the comparisons between arms, which all use one table.

A zero in that table is a tabulated zero. The $\sigma$-profile is built by binning a finite set of
surface segments, so a bin can be empty because no segment of the molecule falls in it as well as
because the molecule cannot carry charge density of that value. The distinction does not affect the
comparisons reported here, all of which treat the tabulation as the reference. It does bound how
strongly the finding may be stated: what is measured is a departure from the tabulated
distribution, not a departure from a proven absence of density.

The profile head reads a role-dependent representation, so a molecule receives two profiles. That
is a choice made in this pipeline, and a reader is entitled to ask whether the defects reported
above are consequences of it. For two of them the question is settled by measurement. The window
displacement of Section~\ref{sec:windows} and the residual population of
Section~\ref{sec:pedestal} are the same on each of the two profiles taken separately, with the
difference between roles no larger than the spread across seeds. The violation of the reference
state in Section~\ref{sec:refstate} is a direct consequence of the choice and would disappear under
a role-independent head, at a cost in accuracy that is not measured here. The non-additivity of
Section~\ref{sec:nonadditive} is a property of the segment-activity calculation rather than of the
head, since it concerns how the closure responds to an internally inconsistent profile. An arm
trained with a role-independent head would settle the remaining cases directly, and has not been
run.

Equation~\eqref{eq:sle} omits the heat-capacity difference between the solid and the subcooled
liquid, as is usual. Where the temperature lies far below the melting point that term is not
negligible, and it adds a further source of indeterminacy to the crystal contribution beyond the
collinearity reported in Section~\ref{sec:ident}.

The activity coefficients used to score the closure are at infinite dilution, while the solubility
the pipeline predicts is a finite mole fraction at saturation. The two are the same functional of
the same profiles evaluated at different compositions. The comparison in
Section~\ref{sec:twosided} therefore reports the direction of the substitution's effect on each,
not a single error carried between them.

This work does not report a more accurate solubility model. On the scaffold-disjoint split used
throughout, the physics path trails a matched direct control sharing the same encoder in the mean
over seeds, and both sit near a random-forest baseline. Both are well behind the leading direct
models on extrapolation to new solutes \cite{osanloo2024}, and absolute errors are within reach of
the disagreement between independent measurements of the same system. No published model was
retrained here, so the comparison with prior work is a comparison of design choices rather than of
performance.

Results are reported over five training seeds, and for some quantities that is not enough. The cost of the spurious density, expressed in the error of the predicted activity coefficients, varies across seeds by several times the spread of the arms themselves. Its sign is therefore not determined by the present evidence. What is determined is that it is not large and positive. The
asymmetry of the cost map between the two sides of the charge-density axis depends on which set of
molecule pairs it is computed over, and is reported for both.

The closure is the 2002 formulation of COSMO-SAC with a single hydrogen-bonding term, run without
the combinatorial contribution. That term needs molar volumes, which the source database carries
but the extracted table used here does not, and which the network could supply with a second scalar
head beside the one that already predicts cavity area. Its omission is a choice made when the table
was extracted rather than a constraint. The term vanishes for a pair of a molecule with itself, so
Section~\ref{sec:refstate} is unaffected; it would change the magnitudes in
Sections~\ref{sec:cost} and~\ref{sec:attribution}, where the size of the effect is the result. Later
formulations split the hydrogen-bonding term by segment type and would change where on the
charge-density axis the cost steps, though not the existence of a step. The segment calculation is
iterated to a fixed count; the count used at training and at scoring is recorded with each run,
because the two need not agree and the predicted activity coefficients depend on them.

Finally, the measurements concern substitution at prediction time on an intermediate that was
trained through the closure. They do not establish what would happen in a model whose intermediate
is supervised externally and frozen. That design is the one the results recommend, and testing it
requires training such a model rather than measuring this one.

\section{Conclusions}
\label{sec:conclusions}

A network trained inside a fixed thermodynamic relation fills the relation's input slot with values
that are not the physical quantity the slot names. The learned $\sigma$-profile places
screening-charge density where the reference tabulation carries none, concentrating in the windows
that trigger the model's hydrogen-bonding term. It differs depending on whether the molecule is treated as solute or solvent, so a molecule is not predicted to be ideal in itself. And it cannot be corrected in parts: replacing a fraction of it with reference values is often worse than replacing none. The output layer is not the cause, and the error that reaches the thermodynamics
lies in the magnitude of the profile where it is populated rather than in the spurious density
where it should be empty.

Supplying the reference profile improves the activity coefficients the profile exists to produce
and degrades the solubility the pipeline predicts. The interpretation that makes physics-informed
models attractive, that a better value for an internal quantity can be looked up and handed to the
model, holds only where that quantity is supervised externally and shielded from the target's
gradient. Where the intermediate is trained end to end, it is an internal coordinate with a
physical name, and the operations the name licenses should be measured before they are used.
